\subsection{Radiation Introduction}
\label{sec: Radiation Background}

Galactic cosmic rays (GCRs) are primary particles produced outside of the solar system \cite{GCRs}. By colliding with other particles, GCRs can induce air showers. These air showers consist of particles, known as secondary particles, which can be detected using sophisticated equipment. By measuring the secondary particles, the primary particle can be characterized its energy and direction. The distribution of secondary particles \cite{Frank} is most dense in the troposphere. In terms of the SORA payload, this means that the bulk of the measured data during ascent and descent will be of secondary particles.

\subsection{MiniPIX Detector}
The MiniPIX \cite{silicon_sensor} is a silicon-based particle detector integrated with a Timepix \cite{timepix} chip. The MiniPIX is a result of the Medipix Collaboration \cite{medipix}. When an energetic particle is incident upon the sensor, energy is deposited in the detector, and the resulting data is recorded. Figure \ref{fig:MiniPIXExample} gives an example of this interaction.

\begin{figure}[h!]
  \centering
  \includegraphics[width=\textwidth]{figures/duck.jpg}
  \caption{look at this damn fine \tiny example output from the silicon-based hybrid pixel detector built by ADVACAM which utilizes the technology developed by CERN's Medipix Collaboration}
  \label{fig:MiniPIXExample}
 \end{figure}
